\clearpage
\SourceLeaf{141}
\begin{center}
[II.]
\end{center}

\begin{center}
AUFGABE.
\end{center}

\bigskip
\hrule
\bigskip

Es sind $p$ Plätze vorhanden, auf welchen $m$ Gegenstände vertheilt werden,
ganz nach Zufall, wobei angenommen wird, für jeden Gegenstand habe jeder
Platz gleiche Wahrscheinlichkeit.

Die Wahrscheinlichkeit, dass die Anzahl der so besetzten Plätze $= m-n$
sein werde, bezeichnen wir mit $(m,n)$ und setzen $\dfrac{1}{p}=x$.


Man hat dann:

\[
\begin{array}{l}
(2,0) = 1-x \\
(2,1) = x \\
(3,0) = (1-x)(1-2x) \\
(3,1) = 3x(1-x) \\
(3,2) = xx \\
(4,0) = (1-x)(1-2x)(1-3x) \\
(4,1) = 6x(1-x)(1-2x) \\
(4,2) = 7xx(1-x) \\
(4,3) = x^3 \\
(5,0) = (1-x)(1-2x)(1-3x)(1-4x) \\
(5,1) = 10x(1-x)(1-2x)(1-3x) \\
(5,2) = 25xx(1-x)(1-2x) \\
(5,3) = 15x^3(1-x) \\
(5,4) = x^4.
\end{array}
\]
\begin{center}
Coefficienten für die folgenden Werthe\\
von $m$:

\smallskip
\begin{tabular}{c||c|c|c|c|c}
$n$ & 5 & 6 & 7 & 8 & 9 \\ \hline\hline
0 & 1 & 1 & 1 & 1 & 1 \\ \hline
1 & 10 & 15 & 21 & 28 & 36 \\ \hline
2 & 25 & 65 & 140 & 266 & 462 \\ \hline
3 & 15 & 90 & 350 & 1050 & 2646 \\ \hline
4 & 1 & 31 & 301 & 1701 & 6951 \\ \hline
5 & & 1 & 63 & 966 & 7770 \\ \hline
6 & & & 1 & 127 & 3025 \\ \hline
7 & & & & 1 & 255 \\ \hline
8 & & & & & 1 \\
\end{tabular}

\end{center}
